% =====================================================================
%  CHAPTER 12: 生命质量评价（对应大纲第十二章）
%  参考PPT：第十二章 生命质量评价
% =====================================================================
\chapter{生命质量评价}
\setcounter{chapter}{12}
\chapterintro{
掌握生命质量（QOL）的概念与构成；掌握生命质量的评价方法（难点*）；熟悉生命质量的测量工具（如SF-36、WHOQOL等）；熟悉生命质量评价的应用（QALY计算等）；了解健康相关生命质量（HRQOL）的概念。
}

\section{生命质量的概念与构成}

\begin{defbox}
\textbf{生命质量（Quality of Life, QOL）：}WHO定义——不同文化和价值体系中的个体对与他们的目标、期望、标准以及所关心的事情有关的生存状态的体验。

\textbf{生命质量的构成：}(1)生理状态（活动能力、角色功能等）；(2)心理状态（情绪、认知、自我评价等）；(3)社会功能状态（社会交往、社会角色、社会支持等）；(4)主观判断与满意度（对自身健康的总体感受和满意度）；(5)对疾病和健康的特殊感受（QOL的疾病特异性维度）。

\textbf{健康相关生命质量（Health-Related Quality of Life, HRQOL）：}指个体对疾病和特定健康状态对其生活质量所造成影响的主观体验和评价，是QOL在健康领域的应用。
\end{defbox}

\section{生命质量的评价方法（难点*）}

\begin{keybox}
\textbf{生命质量的评价内容：}躯体功能（活动能力、日常生活能力、角色功能）、心理功能（情绪、认知、心理应激）、社会功能（社会交往、社会支持）、主观健康（总体健康自评、健康满意度）。

\textbf{评价方法：}
\begin{enumerate}[leftmargin=*]
    \item \imp{访谈法}：通过面对面或电话访谈收集信息
    \item \imp{观察法}：通过直接观察个体行为来评价
    \item \imp{量表法（最常用）}：使用标准化量表进行测量
\end{enumerate}

\textbf{常用测量工具：}SF-36（36条目简明健康调查量表，含13个维度，是国际上最常用的普适性量表之一）、WHOQOL（WHO生命质量量表，含100个问题或简化为26个问题的WHOQOL-BREF）。

\textbf{量表评价要点：}信度（可靠性）、效度（准确性）、反应度（敏感性）。
\end{keybox}

\section{生命质量评价的应用}

\begin{defbox}
\textbf{生命质量评价的应用领域：}
\begin{enumerate}[leftmargin=*]
    \item 临床治疗方案的评价与选择
    \item 人群和患者的健康状况评价
    \item 预防性干预及保健措施的效果评价
    \item 卫生资源配置与利用的决策
    \item 健康影响因素与防治重点的选择
\end{enumerate}

\textbf{QALY（Quality-Adjusted Life Years, 质量调整生存年）：}

QALY = $\sum$ Wk × Yk（Wk为处于k状态的权重值，Yk为处于k状态的年数）。是一种将生存质量和生存时间结合的综合指标，用于比较不同卫生干预措施的成本-效果（COST/QALY）。
\end{defbox}

\section{复习题}

\begin{enumerate}[leftmargin=*, itemsep=8pt]
    \item \textbf{【名词解释】}生命质量（QOL）；健康相关生命质量（HRQOL）；QALY
    \item \textbf{【简答题】}简述生命质量的构成（五个维度）。
    \item \textbf{【简答题】}简述生命质量评价的主要方法和内容。
    \item \textbf{【简答题】}简述生命质量评价的应用领域。
\end{enumerate}

\subsection*{参考答案要点}

\begin{enumerate}[leftmargin=*, itemsep=5pt]
    \item \textbf{QOL}：WHO定义——不同文化和价值体系中的个体对与他们的目标/期望/标准及所关心事情有关的生存状态的体验。\textbf{HRQOL}：个体对疾病和特定健康状态对其生活质量影响的主观体验和评价。\textbf{QALY}：质量调整生存年 = Wk×Yk，将生存质量和生存时间结合的综合指标，用于比较不同卫生干预措施的成本-效果。
    \item 五个维度：(1)生理状态（活动能力/角色功能）；(2)心理状态（情绪/认知/自我评价）；(3)社会功能状态（社会交往/角色/支持）；(4)主观判断与满意度（总体健康感受/满意度）；(5)对疾病和健康的特殊感受（疾病特异性维度）。
    \item 评价内容：躯体功能+心理功能+社会功能+主观健康。评价方法：访谈法、观察法、量表法（最常用，如SF-36、WHOQOL）。评价要点：信度、效度、反应度。
    \item (1)临床治疗方案评价与选择；(2)人群和患者健康状况评价；(3)预防性干预及保健措施效果评价；(4)卫生资源配置与利用决策；(5)健康影响因素与防治重点选择。
\end{enumerate}

\newpage
